\documentclass[11pt]{article}
\usepackage[margin=2.3cm]{geometry}

\usepackage{graphicx}
\usepackage{float}
\usepackage{amsmath,amssymb,amsfonts}
\usepackage{hyperref}
\usepackage{booktabs}
\usepackage{xcolor}
\usepackage{array}

\hypersetup{
  colorlinks=true,
  linkcolor=blue!50!black,
  citecolor=green!50!black,
  urlcolor=blue!50!black,
}

\newcommand{\su}{\mathfrak{su}}
\newcommand{\half}{\tfrac{1}{2}}
\newcommand{\CA}{C_A}
\newcommand{\ksm}{\kappa_\text{SM}}

\begin{document}

\title{One Knot, All Forces:\\[4pt]
The Gravitational Constant from the\\
Poincar\'e Sphere of the Trefoil\\[2pt]
\normalsize (Paper 15)}

\author{
  James~P.~Galasyn$^{1}$ \quad and \quad Claude~Th\'eodore$^{2}$\\[2pt]
  \small $^1$Independent researcher \\
  \small $^2$Anthropic
}

\date{\today}

\maketitle

\begin{abstract}
Papers 13--14 derived eighty Standard Model observables from the
trefoil knot $T(2, 3)$.  This paper argues that $(+1)$-Dehn surgery
on the trefoil---which produces the Poincar\'e homology sphere
$S^3/2I$ and the icosahedral completion $A_5$ used for nuclear
binding in the companion Paper~17---also organises the integers
appearing in the gravitational constant.

The first $2I$-invariant Laplacian eigenvalue on $S^3/2I$ is
$\lambda_1 = 168$, which factorises two complementary ways:
$168 = 8 \times 21$ (matching the spinor and adjoint dimensions
of $\mathrm{Spin}(7) = B_3$) and $168 = 7 \times 24 =
|\mathrm{Irr}(2T)| \times |2T|$ via the McKay correspondence on
the binary tetrahedral group.  The integer $7$ appears
simultaneously as the spectral quotient $(n_1{+}2)/2$, as the
McKay-node count on the extended $E_6$ diagram, and as the
vector rep dimension of $\mathrm{Spin}(7)$; $2T$ embeds explicitly
into $\mathrm{Spin}(7)$ via the octonion Clifford algebra
$\mathrm{Cl}(0, 7)$ (\S\ref{sec:spectrum}, \S\ref{sec:mechanism}).
These integers enter the proposed hierarchy
\begin{equation*}
  \frac{m_e}{m_\text{Pl}}
  = \frac{8}{7}\left(1 + \frac{\alpha}{7}\right)
    \alpha^{21/2} ,
\end{equation*}
where the exponent arises, under $PSL(2, 7)$-equivariance, from
a Wilson-line amplitude on the $21$-edge complete graph $K_7$
(the Heawood map on the Heegaard torus of $S^3/2I$; $21 =
\dim\mathfrak{so}(7)$).  The resulting gravitational constant
$G = (8/7)^2(1{+}\alpha/7)^2\,\alpha^{21}\,\hbar c/m_e^2$ agrees
numerically with CODATA at $\mathbf{0.013\%}$ using the
NWT-derived $\alpha$, or $0.029\%$ using the measured $\alpha$.

Within the present accounting, no integer in the formula is
introduced as an independent input beyond those already
generated by the trefoil's surgery and spectrum; the cinquefoil
$T(2, 5)$ and heptafoil $T(2, 7)$ appear as derived labels rather
than additional free knots.  The theory still contains a
distinguished mass scale $v_\text{EW}$ (equivalently $m_e$),
which we treat as a reference scale setting the unit of
measurement; the zero-parameter claim refers to the
\emph{dimensionless} content.

Reproducibility: all analysis code (Laplacian eigenvalues on
$S^3/2I$, the $\mathrm{Cl}(0, 7)$ construction of $2T \subset
\mathrm{Spin}(7)$, the one-loop BPS pipeline, the $SU(N)$
robustness scan, and the $G$-formula verification) is available
at \url{https://github.com/JimGalasyn/null-worldtube}.
\end{abstract}

\tableofcontents

%=====================================================================
\section{Introduction: The Last Ratio}
\label{sec:intro}
%=====================================================================

The Standard Model of particle physics, combined with general
relativity, contains approximately twenty-one free parameters:
twenty dimensionless quantities (coupling constants, mass ratios,
mixing angles) and one overall mass scale.  Papers~13--14 of
this series~\cite{paper13,paper14} derived all twenty
dimensionless SM parameters from the topology of the trefoil
knot $T(2,3)$.  The companion Paper~17~\cite{paper17} extends
the framework to nuclear physics, deriving the binding energy
coefficients from the same topology.

One dimensionless ratio remained: $m_e/m_\text{Pl} \approx
4.2 \times 10^{-23}$, the ratio of the electron mass to the
Planck mass.  This is the \emph{hierarchy problem}: why is
gravity $10^{42}$ times weaker than electromagnetism?

This paper eliminates the hierarchy as a free parameter.  The
ratio $m_e/m_\text{Pl}$ is determined by the \emph{spectral
geometry of the Poincar\'e homology sphere} $S^3/2I$---the
manifold obtained by $(+1)$-Dehn surgery on the trefoil, the
same surgery that produced the icosahedral group $A_5$ in
Paper~17.

The result is a \textbf{one-knot accounting} in which the
trefoil organises the dimensionless content of the framework
to within the precision of the present calculations.


%=====================================================================
\section{From the Trefoil to the Poincar\'e Spectrum}
\label{sec:spectrum}
%=====================================================================

\subsection{Surgery and the Poincar\'e sphere}

$(+1)$-Dehn surgery on the left-handed trefoil produces the
Poincar\'e homology sphere
$\Sigma(2,3,5) = S^3/2I$~\cite{Rolfsen}, where $2I$ is the
binary icosahedral group of order~120.  This is a standard
result in low-dimensional topology, used in Paper~14 to
derive $k_\text{sat} = |\mathrm{Irr}(A_5)| = 5$ and in
Paper~17 to derive the nuclear SEMF coefficients.

\subsection{The Laplacian on $S^3/2I$}

The eigenvalues of $-\nabla^2$ on $S^3$ of radius $R$ are
$\lambda_n = n(n+2)/R^2$ for $n = 0, 1, 2, \ldots$, with
degeneracy $(n+1)^2$.  On the quotient $S^3/2I$, only the
$2I$-invariant harmonics survive: those for which the
spin-$n/2$ representation of $SU(2)$, restricted to $2I$,
contains the trivial representation.

The trivial-representation multiplicity is
\begin{equation}
  m(n) = \frac{1}{|2I|}\sum_{g \in 2I} \chi_{n/2}(g)\,,
\end{equation}
where $\chi_j(g) = \sin((2j{+}1)\theta_g/2)/\sin(\theta_g/2)$
is the spin-$j$ character at rotation angle~$\theta_g$, summed
over the nine conjugacy classes of~$2I$.

The result: $m(n) = 0$ for $1 \leq n \leq 11$, and
\begin{equation}
\label{eq:lambda1}
\boxed{
  n_1 = 12\,,\qquad
  \lambda_1 = 12 \times 14 = 168\,.
}
\end{equation}
The integer $n_1 = 12$ is the number of vertices of the
icosahedron---the defining geometry of~$A_5$.  The gap between
the constant mode ($n = 0$) and the first excitation ($n = 12$)
is the \emph{spectral gap} of the Poincar\'e sphere.

\subsection{Factorisation of $\lambda_1$}

The eigenvalue $168$ factorises as
\begin{equation}
  168 = 8 \times 21 = (\CA + 1) \times
  \dim(\wedge^2\,\mathbf{C}^{\CA})\,,
  \qquad \CA = \frac{n_1 + 2}{2} = 7\,,
\end{equation}
where $\wedge^2 \mathbf{C}^{\CA}$ is the antisymmetric square of
the defining (fundamental) representation of $\su(\CA)$, with
dimension $\CA(\CA{-}1)/2$.  (The ambient algebra $\su(\CA)$
itself has dimension $\CA^2 - 1 = 48$; we are \emph{not}
referring to $\wedge^2$ of the adjoint.)
The integer $\CA = 7$ emerges from the spectral data
($n_1 = 12 \Rightarrow n_1 + 2 = 14 = 2\CA$), not from a
separate knot.  The factorisation is unique in the following
precise sense: $8 \times 21$ is the only splitting of $168$
into $(\CA{+}1)$ and $\CA(\CA{-}1)/2$ for any integer $\CA$.
Proof: setting $(\CA{+}1) \cdot \CA(\CA{-}1)/2 = 168$ gives
$\CA^3 - \CA = 336$, a monotone-increasing cubic on
$\CA \geq 1$ with unique positive integer root $\CA = 7$
(since $7^3 - 7 = 336$).  This means the spectral datum
$\lambda_1 = 168$ singles out $\CA = 7$ with no alternative.

\medskip
\noindent\textbf{Group-theoretic origin via the McKay correspondence.}
The factorisation $168 = 8 \times 21$ is not merely arithmetic but
group-theoretic.  The McKay correspondence pairs finite subgroups
of $SU(2)$ with extended Dynkin diagrams; for the binary tetrahedral
group $2T$ (order 24) the corresponding diagram is extended $E_6$,
which has exactly $7$ nodes corresponding to the $7$ irreducible
representations of $2T$ with dimensions $(1,1,1,2,2,2,3)$.  Thus
\begin{equation}
  \CA \;=\; 7 \;=\; |\mathrm{Irr}(2T)|,\qquad
  168 \;=\; 7 \times 24 \;=\; |\mathrm{Irr}(2T)| \times |2T|.
\end{equation}

These are not two different factorisations of two different
numbers: $168 = 8 \times 21 = 7 \times 24$ is the \emph{same}
integer viewed two complementary ways.  The factorisation
$8 \times 21 = (7{+}1) \cdot 7(7{-}1)/2$ packages $\lambda_1$ in
terms of the Spin$(7) = B_3$ representation dimensions
$(\text{vector}, \text{spinor}, \text{adjoint}) = (7, 8, 21)$; the
factorisation $7 \times 24 = |\mathrm{Irr}(2T)| \cdot |2T|$ packages
the same $\lambda_1$ in terms of the McKay correspondence on the
binary tetrahedral group.  The two views share the integer~$\CA = 7$:
spectrally it is $(n_1{+}2)/2$, group-theoretically it is the number
of nodes of the extended $E_6$ Dynkin diagram, and representationally
it is $\dim(\text{vector rep of } \mathrm{Spin}(7))$.  All four
integers $(7, 8, 21, 24)$ play distinct, named roles in the
gravitational hierarchy formula and the trefoil $2I$-orbit
combinatorics; the ``uniqueness'' claim above refers specifically
to the $(\CA{+}1) \times \CA(\CA{-}1)/2$ form, which singles out
$\CA = 7$.

The group $2T$ arises naturally in the setup: the trefoil's
stabiliser in $2I = \pi_1(S^3/2I)$ is $\mathbb{Z}_5$, and since
$\gcd(24,5)=1$ the decomposition $2I = 2T \cdot \mathbb{Z}_5$ is
bijective as a set, so the orbit of the trefoil under $2I$ has
exactly $24$ elements, in natural bijection with~$2T$.

\subsection{Robustness: exclusion of alternatives and
look-elsewhere correction}
\label{sec:robustness}

A numerical coincidence at the $0.03\%$ level demands scrutiny.
We address this in two stages: first the unconstrained search
space, then the structurally constrained one.

\medskip
\noindent\textbf{Unconstrained search.}
Among $\sim\!4000$ candidate formulas of the form
(simple fraction) $\times \alpha^{n/2}$ with
$p, q \leq 15$ and $n \in \{18, \ldots, 25\}$, the probability
that a random target in the relevant range is matched at
$\lesssim 0.03\%$ precision by \emph{some} candidate is
$\sim 9\%$.  Taken at face value, this is a look-elsewhere
effect large enough that bare numerical agreement alone would
\emph{not} be compelling evidence.

\medskip
\noindent\textbf{Structural constraints cut the search space to
$\mathcal{O}(1)$.}
The formula
$m_e/m_\text{Pl} = (\CA{+}1)/\CA \times (1 + \alpha/\CA) \times
\alpha^{\CA(\CA-1)/4}$
is not a free parametrisation of the $4000$-formula family: it is
the \emph{only} formula in that family admitting the $(\CA{+}1)
\times \CA(\CA{-}1)/2$ factorisation of $\lambda_1 = 168$ (the
uniqueness statement at the end of \S\ref{sec:spectrum}).  Once
we require the integers to arise from the spectral datum
$\lambda_1 = n_1(n_1{+}2) = 168$ on $S^3/2I$ via that factorisation
rule, the search space collapses to a single $\CA$-family with
$\CA$ determined by the spectral quotient $(n_1{+}2)/2$.  The
$SU(N)$ scan spans $N \in \{2, \ldots, 14\}$---$13$ candidates
rather than $4000$:

\begin{itemize}\setlength\itemsep{1pt}
\item Only $\CA = 7$ matches within $1\%$; the nearest competitor
$SU(6)$ fails by a factor of $10^8$ (the exponential dependence
on $\dim(\wedge^2) = \CA(\CA{-}1)/2$ is steep).
\item Within $\su(7)$, only $\dim(\wedge^2) = 21$ works;
$\dim(\text{adj}) = 48$, $\dim(\text{fund}) = 7$, and
$\dim(\text{Sym}^2) = 28$ each fail by $\gtrsim 10^7$.
\end{itemize}

A Bonferroni-style bound on the post-selection probability is
therefore
\begin{equation*}
  P\bigl(\text{random match at } 0.03\%\ \big|\ \text{structural family}\bigr)
  \;\lesssim\; 13 \times 2 \times 10^{-4}
  \;\approx\; 3 \times 10^{-3},
\end{equation*}
rather than the unconstrained $9 \times 10^{-2}$---a factor of
$\sim 30$ reduction.  In practice the reduction is much
stronger because the $\sim 13$ structural candidates are not
uniformly distributed in the target range: only $\CA = 7$ falls
within $1\%$, so the \emph{effective} number of same-precision
competitors is closer to $1$, and the conditional $p$-value is
correspondingly $\sim\!10^{-4}$.

\medskip
\noindent\textbf{The spectral origin is the key.}
The statistical argument depends crucially on the integers
$(7, 8, 21)$ being fixed by the spectral datum before the
numerical fit, not after.  The chain
\begin{center}
trefoil $\xrightarrow{(+1)\text{-Dehn}}$ $S^3/2I$
$\xrightarrow{\text{Laplacian}}$ $\lambda_1 = 168$
$\xrightarrow{\text{McKay/factorisation}}$ $\CA = 7$
\end{center}
produces $\CA = 7$ \emph{without any reference to $m_e/m_\text{Pl}$}.
It is this external derivation of the integers that lets the
$0.013\%$ CODATA agreement be read as evidence rather than
curve-fitting.  The reproducibility pipeline for each link
(Laplacian eigenvalue computation, McKay decomposition,
Cl$(0,7)$ construction, $SU(N)$ scan) is contained in the code
repository cited in the Code availability note at the end of
this paper.


%=====================================================================
\section{The Gravitational Hierarchy}
\label{sec:hierarchy}
%=====================================================================

\subsection{The formula}

We propose, as an ansatz motivated by the structural chain above,
that the electron-to-Planck mass ratio takes the form
\begin{equation}
\label{eq:hierarchy}
\boxed{
  \frac{m_e}{m_\text{Pl}}
  = \frac{\CA + 1}{\CA}
    \left(1 + \frac{\alpha}{\CA}\right)
    \alpha^{\dim(\wedge^2)/2}\,,
  \qquad \CA = 7\,,
}
\end{equation}
where the integers are chosen to align with the Poincar\'e
spectrum:
\begin{itemize}\setlength\itemsep{0pt}
  \item $\CA = 7 = (n_1 + 2)/2$ from the icosahedral harmonic
    level $n_1 = 12$.
  \item $\dim(\wedge^2) = 21 = \lambda_1/(\CA + 1) = 168/8$
    from the eigenvalue factorisation.
  \item $(\CA + 1)/\CA = 8/7$ from the same factorisation.
  \item $\alpha = 1/(25\pi\sqrt{3}+1)$ from the trefoil's AB
    holonomy (Papers~13--14).
\end{itemize}

\subsection{NLO correction: the recurring $(1 + \alpha/\CA)$ pattern}
\label{sec:nlo}

We interpret the factor $(1 + \alpha/7)$ as a one-loop-like
correction associated with the $\CA = 7$ sector.  The same
$(1 + \alpha/\CA)$ pattern recurs in the other NWT examples
examined so far:

\begin{table}[h]
\centering
\caption{Universal $(1 + \alpha/\CA)$ NLO pattern.}
\label{tab:nlo}
\begin{tabular}{llcl}
\toprule
Quantity & Group & NLO factor & Residual \\
\midrule
$\alpha_s(M_Z)$ & $\su(3)$ & $(1 + \alpha/3)$ & $0.08\%$ \\
SEMF $E_1$ & $\su(3)$ & $(1 - 9\alpha)$ & $1.66\%$ RMS \\
$m_e/m_\text{Pl}$ & $\su(7)$ & $(1 + \alpha/7)$ & $\mathbf{0.006\%}$ \\
\bottomrule
\end{tabular}
\end{table}

\subsection{The gravitational constant}

\begin{equation}
\label{eq:G}
\boxed{
  G = \left(\frac{8}{7}\right)^{\!2}
  \left(1 + \frac{\alpha}{7}\right)^{\!2}
  \alpha^{21}\,\frac{\hbar c}{m_e^2}\,.
}
\end{equation}

\begin{table}[h]
\centering
\caption{Predicted $G$ at leading order and NLO.}
\label{tab:G}
\begin{tabular}{lccc}
\toprule
 & $G$ [$10^{-11}$~m$^3$\,kg$^{-1}$\,s$^{-2}$]
 & NWT $\alpha$ & CODATA $\alpha$ \\
\midrule
LO: $(8/7)^2\,\alpha^{21}$ & $6.660$ & $-0.22\%$ & $-0.24\%$ \\
NLO: $+\,(1{+}\alpha/7)^2$ & $6.673$ & $\mathbf{-0.013\%}$
  & $-0.029\%$ \\
CODATA 2022 & $6.674$ & --- & --- \\
\bottomrule
\end{tabular}
\end{table}

\noindent With the NWT-derived $\alpha = 1/(25\pi\sqrt{3}+1)$,
the NLO prediction matches CODATA at $0.013\%$; with the
CODATA-measured $\alpha$, the match is $0.029\%$.  Both are
within the spread of independent $G$ measurements ($\sim 0.05\%$).
The factor-of-two difference arises because
$G \propto (m_e/m_\text{Pl})^{-2}$, doubling any residual in
the mass ratio.


%=====================================================================
\section{One Knot, Not Three}
\label{sec:one-knot}
%=====================================================================

Papers~13--14 used three torus knots: the trefoil $T(2,3)$, the
cinquefoil $T(2,5)$, and (implicitly) the heptafoil $T(2,7)$.
A natural reading was ``three knots determine three forces.''
The present accounting suggests a tighter bookkeeping: the
cinquefoil and heptafoil can be viewed as derived organizational
labels within a trefoil-based scheme, rather than as additional
free topological inputs.

\begin{table}[h]
\centering
\caption{All integers trace to the trefoil, via three routes.}
\label{tab:one-knot}
\begin{tabular}{rclll}
\toprule
Integer & Value & Identity & Route & Derivation \\
\midrule
\multicolumn{5}{l}{\textit{Direct (crossing algebra):}}\\
$\CA(\su(3))$ & 3 & trefoil crossings & direct & Paper~14 \\
\midrule
\multicolumn{5}{l}{\textit{Dehn surgery (trefoil $\to$ $A_5$):}}\\
$|\mathrm{Irr}(A_5)|$ & 5 & cinquefoil winding
  & surgery & Paper~14 \\
\midrule
\multicolumn{5}{l}{\textit{Spectral geometry (Poincar\'e sphere $S^3/2I$):}}\\
$(n_1{+}2)/2$ & 7 & heptafoil Casimir
  & spectral & this paper \\
$\lambda_1/\dim(\wedge^2)$ & 8 & correction factor
  & spectral & this paper \\
$\lambda_1/(\CA{+}1)$ & 21 & hierarchy exponent
  & spectral & this paper \\
$\lambda_1(S^3/2I)$ & 168 & spectral gap
  & spectral & this paper \\
\bottomrule
\end{tabular}
\end{table}

The cinquefoil $T(2,5)$ carries the integer
$5 = |\mathrm{Irr}(A_5)|$, which was derived from trefoil
surgery in Paper~14.  The heptafoil $T(2,7)$ carries the
integers $7$, $8$, and $21$, which align with the Poincar\'e
spectrum $\lambda_1 = 168$ on $S^3/2I$.  Within the present
accounting, neither knot is introduced as an additional free
topological input; their integer content is encoded in the
trefoil / surgery / spectrum chain.

The remaining non-trefoil ingredient is $\CA(\su(2)) = 2$ from
the Hopf link---the linking topology that binds trefoil carriers
into nuclei (Paper~17) and generates the electroweak $\su(2)$.


%=====================================================================
\section{Cosmological Phase Transitions}
\label{sec:cosmology}
%=====================================================================

Though all integers trace to the trefoil, they play
\emph{three distinct physical roles}---not three copies of the
same role.  The following cosmological interpretation is
\emph{speculative}: it provides a physical narrative for the
mathematical structure but does not constitute a dynamical
derivation from evolution equations.

\begin{table}[h]
\centering
\caption{Three roles, one knot.}
\label{tab:roles}
\begin{tabular}{llll}
\toprule
Role & Integer source & What it determines & Epoch \\
\midrule
\textbf{Gauge} & crossings ($q{=}3$)
  & colour force, $\alpha$ & active today \\
\textbf{Architectural} & surgery ($q{=}5$)
  & gauge embedding, $k_\text{sat}$ & frozen (GUT) \\
\textbf{Structural} & spectrum ($\lambda_1$)
  & gravitational hierarchy $G$ & imprinted (Planck) \\
\bottomrule
\end{tabular}
\end{table}

As the universe cools from the Planck temperature, these roles
freeze out in reverse order:
\begin{enumerate}\setlength\itemsep{3pt}
  \item \textbf{Planck epoch} ($T \sim 10^{32}$~K):
    the Poincar\'e spectrum is established.  $G$ crystallises
    from $\lambda_1 = 168$.

  \item \textbf{GUT epoch} ($T \sim 10^{28}$~K):
    $A_5$ symmetry breaks.
    $\su(5) \to \su(3) \times \su(2) \times \mathfrak{u}(1)$.
    The gauge structure and $\alpha$ are fixed.

  \item \textbf{QCD epoch} ($T \sim 10^{12}$~K):
    trefoil vortices (baryons) condense.  The strong force
    confines.  This is the world we observe.
\end{enumerate}

Each frozen epoch leaves a relic:
$G$ is a \emph{spectral fossil} of the Poincar\'e sphere;
$\alpha$ and the gauge group are \emph{surgical fossils}
of $A_5$;
baryons and confinement are the \emph{living topology} of the
trefoil.


%=====================================================================
\section{Inertia, Gravity, and the Equivalence Principle}
\label{sec:equivalence}
%=====================================================================

\textbf{Inertia} is the condensate's resistance to vortex
displacement: accelerating a vortex ring requires reorganising
the surrounding condensate.  The BPS energy $\mu \times 2\pi R$
is the inertial mass.

\textbf{Gravity} is the condensate's resistance to curvature:
$G \propto 1/v_\text{EW}^2$ (Eq.~\ref{eq:G}), so a stiffer
condensate gives weaker gravity.

\textbf{The equivalence principle} follows: both are responses
of the same condensate.  There is no separate gravitational
charge.  Inertial mass equals gravitational mass because they
are the same condensate property measured in two ways.

\medskip
\noindent\textbf{Status of this section.}
The arguments above are \emph{heuristic}: they are intended to
illustrate how the equivalence principle would follow from the
condensate picture, not to derive it rigorously from the
Lagrangian.  A first-principles derivation requires showing that
the effective graviton propagator and the effective inertial
response on $S^3/2I$ both emerge from the same BPS sector of
$\mathcal{L}_2$ (\S\ref{sec:lagrangian-context}) via the Volovik
analog-gravity mechanism; this is part of the open dynamical
programme of the companion Paper~16.


%=====================================================================
\section{A Zero-Parameter Theory (for dimensionless content)}
\label{sec:zero}
%=====================================================================

With the gravitational hierarchy in hand, the dimensionless
parameter count is:

\begin{table}[h]
\centering
\caption{From twenty parameters to zero.}
\label{tab:params}
\begin{tabular}{lcc}
\toprule
 & SM + GR & NWT \\
\midrule
Dimensionless parameters & $\sim 20$ & 0 \\
Mass scale & 1 ($v_\text{EW}$) & 1 ($v_\text{EW}$) \\
Gauge group (input) & $SU(3){\times}SU(2){\times}U(1)$ & --- \\
Gauge group (output)$^\dagger$ & --- &
  $\su(3){\times}\su(2){\times}\mathfrak{u}(1)$ \\
$G$ (input) & yes & --- \\
$G$ (output) & --- & yes ($0.013\%$) \\
\bottomrule
\end{tabular}\\[4pt]
{\footnotesize $^\dagger$The script font $\su(n)$ denotes the
Lie \emph{algebra} (generators and commutation relations),
derived from crossing algebra.  The roman $SU(n)$ denotes the
Lie \emph{group}.  The algebra uniquely determines the group
for the SM representations.}
\end{table}

The remaining scale $v_\text{EW}$ is treated here as a
\textbf{reference scale} that sets the unit of measurement.  In
any system of units, one mass must be chosen as the reference;
$v_\text{EW}$ plays this role in the present framework.  Within
that convention, every \emph{dimensionless} ratio, coupling, and
mixing angle is determined by the trefoil's topology to within
the precision of the current calculations.

This is analogous to QCD, where $\Lambda_\text{QCD}$ is
generated by dimensional transmutation rather than being a free
parameter.  Nobody counts $\Lambda_\text{QCD}$ among the ``free
parameters of QCD.''  Similarly, $v_\text{EW}$ is the NWT
condensate's transmutation scale, not a tunable input.

A philosophical objection can be raised: $v_\text{EW}$ still
\emph{does work} in the framework---it sets the absolute scale
of all predictions.  This is correct, but it is equally true
of the kilogram, the second, and the metre.  The distinction is
between a parameter that \emph{could have been different}
(changing $\alpha$ would change the physics) and a scale that
\emph{defines the units} (changing $v_\text{EW}$ changes only
the numbers we write down, not the physical content).  The
$\sim 20$ SM parameters are of the first type; $v_\text{EW}$ is
of the second.  We therefore maintain the ``zero dimensionless
free parameters'' characterisation while acknowledging that one
reference scale remains.


%=====================================================================
\section{Discussion}
\label{sec:discussion}
%=====================================================================

\subsection{What is derived and what is assumed}

The framework assumes:
\begin{itemize}\setlength\itemsep{0pt}
  \item The abelian Higgs Lagrangian at the BPS point.
  \item The trefoil $T(2,3)$ as the fundamental knot.
  \item $3+1$ spacetime dimensions.
  \item $\hbar$ and $c$ as fundamental constants.
\end{itemize}
Everything else is derived:
$\alpha$, $G$, the gauge group, the mass spectrum, the mixing
matrices, the nuclear binding energies.

\subsection{Structural analysis of the exponent}
\label{sec:mechanism}

The formula $m_e/m_\text{Pl} = (8/7)(1{+}\alpha/7)\,
\alpha^{21/2}$ matches CODATA at $0.03\%$, is structurally
unique (\S\ref{sec:hierarchy}), and has every integer traceable
to the Poincar\'e spectrum (\S\ref{sec:spectrum}).  We do not
yet possess a \emph{dynamical derivation} of the formula---an
explicit computation showing that the gravitating BPS theory on
$S^3/2I$ produces $\alpha^{21}$ in the effective coupling.
What we can exhibit is the \emph{structural anatomy} of the
integers.

\medskip
\noindent\textbf{The eigenvalue decomposition.}\\
The spectral gap mode at $n_1 = 12$ has spin $j = 6$ and
Casimir $C_2 = j(j+1) = 42$.  The eigenvalue is
$\lambda_1 = n_1(n_1{+}2) = 168 = 4 C_2$.  The exponent in
the hierarchy formula satisfies
\begin{equation}
\label{eq:exponent-origin}
  21 = \frac{\lambda_1}{\CA + 1} = \frac{C_2}{2}
     = \frac{j(j+1)}{2}
     = \dim(\wedge^2(\CA))\,,
  \qquad \CA = j{+}1 = 7\,.
\end{equation}
The identity $\CA(\CA{-}1)/2 = j(j{+}1)/2$ (forced by
$\CA = j{+}1$) guarantees that $\lambda_1/(\CA{+}1)$
simultaneously equals half the Casimir \emph{and} the
antisymmetric-square dimension.  This is an algebraic fact,
not a dynamical derivation.

\medskip
\noindent\textbf{The antisymmetric structure.}\\
The exponent is $21 = \dim(\wedge^2)$, not $28 = \dim(\mathrm{Sym}^2)$.
This distinction has a physical correlate: the BPS gauge field
$F_{\mu\nu}$ is a 2-form (antisymmetric), so the bilinear
stress-energy $T \sim F \wedge F$ lives in $\wedge^2$ rather
than $\mathrm{Sym}^2$.  If the gauge field were a scalar, the
natural exponent would be $28$ and the prediction would fail.
The correct exponent \emph{requires} the gauge field to be a
1-form---which it is, by construction, in the abelian Higgs
model.  This is the section's strongest structural insight: the
antisymmetric nature of $F$ is not a choice but a consequence
of the Lagrangian, and it selects $21$ over~$28$.

\medskip
\noindent\textbf{What is derived and what is conjectured.}\\
\emph{Derived:} $\lambda_1 = 168$ (spectral geometry);
$\CA = 7 = |\mathrm{Irr}(2T)|$ (McKay correspondence, extended
$E_6$); $|2T| = 24$ (trefoil $2I$-orbit size); $21 = \dim\mathfrak{so}(7)$
(algebraic identity); $21$ not $28$ (antisymmetric $F$);
explicit $2T \subset \mathrm{Spin}(7)$ matrix embedding via the
$\mathrm{Cl}(0, 7)$ octonion Clifford algebra (see
\S\ref{sec:mechanism} below); the $\alpha^{21/2}$ exponent as a
\emph{structural} consequence of the Wilson-line amplitude through
the $21$ $\mathfrak{so}(7)$ generators with $PSL(2,7)$-equivariance
(derivation sketched below).

\emph{Conjectured:} the overall amplitude form.  Specifically:
the $\alpha$-\emph{exponent} $21/2$ is structurally determined
from the $\mathrm{Cl}(0,7)$ / $\mathfrak{so}(7)$ / $PSL(2,7)$
chain (derivation below), and the $(8/7)$ prefactor has a concrete
$\mathrm{Spin}(7)$-invariant source:
\begin{equation}
  \frac{8}{7}
  \;=\; \frac{\dim(\text{spinor})}{\dim(\text{vector})}
  \;=\; \frac{C_2(\text{vector})}{C_2(\text{spinor})}
  \quad\text{for }\mathrm{Spin}(7) = B_3\,.
\end{equation}
These two ratios are the same identity: both the vector and spinor
representations of $B_3$ have Dynkin index $T(\pi) = 2$, so
$\dim(\pi)\,C_2(\pi)$ is equal across them, forcing
$\dim(\text{spinor})/\dim(\text{vector}) = C_2(\text{vector})/C_2(\text{spinor})$.
Physically, this is the ratio of matter-propagator dimensionality
($8$-dim spinor of $\mathrm{Spin}(7)$, realised as the full octonion
space via $\mathrm{Cl}(0, 7)$) to gauge-exchange dimensionality
($7$-dim vector rep, $\operatorname{Im}(\mathbb{O})$) in the
Wilson-line amplitude.

The NLO factor $(1 + \alpha/\CA)$ is the recurring Paper~15
pattern (\S\ref{sec:nlo}, Table~\ref{tab:nlo}) and is consistent
with the generic structure of one-loop gauge-theory corrections
with a $\CA$-dim vector internal line.  One natural realisation
in the $K_7$ / Wilson-line picture: each of the $7$ vertices of
$K_7$ receives a self-energy correction of order $\alpha/\CA^2 =
\alpha/49$, and summing over $\CA = 7$ vertices yields
$7 \cdot \alpha/49 = \alpha/\CA$.  A first-principles derivation
from a specific Lagrangian remains open, but the coefficient
$1/\CA$ is structurally consistent with the $\mathrm{Cl}(0, 7)$
/ $\mathrm{Spin}(7)$ chain and with the NLO pattern recurring
across Paper~15's three examples ($SU(3)$, $SU(3)$,
$\mathrm{Spin}(7)$); the three residuals span $0.006\%$ to
$1.66\%$, so the evidence is suggestive rather than conclusive,
and the SEMF $E_1$ outlier is a target for a proper derivation.

A dynamical derivation would require computing the
\textbf{scalar-sector} self-energy on $S^3/2I$ in the presence
of the BPS gauge field, and showing that the propagator pole at
$\lambda_1$ generates exactly $\alpha^{21}$ in the effective
coupling.  (The pole is unambiguously a scalar-mode pole, not a
TT-graviton pole: on $S^3/2I$ the transverse-traceless tensor
sector at $\lambda = 168$ decomposes under
$SU(2)_L \times SU(2)_R$ as $(j_L, j_R) = (7,5) \oplus (5,7)$,
both of which have \emph{zero} $2I$-invariant subspace.  The
gravitational response at this eigenvalue is therefore carried
by the condensate's compressibility mode, consistent with
Volovik-style analog gravity where the condensate's compressional
response dominates long-wavelength gravitational dynamics.)

The one-loop machinery required for such a dynamical derivation has
been validated on the flat-space $2D$ benchmark.  A full
Alonso-Izquierdo--style Casimir computation for the critically-coupled
Nielsen--Olesen vortex~\cite{AlonsoIzquierdo2016,BaackeLavrelashvili2008},
with the boson, ghost, and Faddeev--Popov sectors treated via heat-kernel
/ Jorgenson--Lang $\zeta$-regularisation, reproduces the literature value
$\Delta\mu \approx -0.279$ in magnitude and sign when the ghost is
assigned its natural two-real-DOF Grassmann normalisation (coefficient
$-\hbar m$ per complex ghost, not $-\hbar m / 2$).  This validates the
pipeline (eigensolver, Seeley coefficients, $\zeta(-1/2)$ extraction,
ghost treatment) that would be re-deployed for the $S^3/2I$ curved-space
computation.

Alternatively, Paper~13's rule of $\sqrt{\alpha}$ per crossing
on \emph{open} (transition) paths connects to the hierarchy via
a specific combinatorial target.  The integer~$21$ is the edge
count of the complete graph $K_7$ on the $7$ McKay nodes of $2T$,
which embeds on the \textbf{Heegaard torus} of $S^3/2I$ as the
Heawood map ($V=7$, $E=21$, $F=14$, Euler characteristic $0$).
$K_7$ admits Eulerian circuits (every vertex has even
degree~$6$), and any such circuit traverses all $21$ edges
exactly once.  Under Paper~13's open-path rule, the product over
$21$~crossings gives exactly $(\sqrt{\alpha})^{21} = \alpha^{21/2}$,
matching the exponent in $m_e/m_\text{Pl}$.  This promotes the
``21~crossings'' conjecture from mode-counting to a specific
diagrammatic target: the $21$-edge $K_7$ Eulerian circuit on the
Heegaard torus.

\medskip
\noindent\textbf{Why this diagram?  The Wilson-line / $\mathfrak{so}(7)$
bridge.}
The $K_7$ Eulerian-circuit prescription is not an ad hoc choice:
it follows from the explicit $\mathrm{Cl}(0, 7)$ octonion Clifford
construction of $2T \subset \mathrm{Spin}(7)$ (\S\ref{sec:mechanism})
together with $PSL(2, 7)$-equivariance.  The $21$ bivectors
$B_{ij} := L_i L_j$ for $1 \leq i < j \leq 7$ span
$\mathfrak{so}(7)$ (rank $21$, Lie-closure verified numerically to
machine precision), and are in natural bijection with the $21$
edges $\{i, j\}$ of $K_7$:
\begin{equation}
  \text{$K_7$ edge $\{i, j\}$}
  \;\longleftrightarrow\;
  \text{$\mathfrak{so}(7)$ generator } B_{ij} = L_i L_j .
\end{equation}
In the Wilson-line expansion of the transition amplitude
$\langle\sigma |\, \hat{O} \, | e\rangle$, each insertion of an
$\mathfrak{so}(7)$ gauge boson contributes one $\sqrt{\alpha}$
factor by Paper~13's open-path rule.  $PSL(2,7)$ acts
$2$-transitively on the $7$ McKay nodes, hence transitively on
the $21$ edges of $K_7$ and on the $21$ generators of
$\mathfrak{so}(7)$.  Any $PSL(2,7)$-equivariant amplitude must
therefore couple to \emph{all} $21$ generators---an amplitude
missing any single generator cannot lie in a single $PSL(2, 7)$
orbit.  The minimal closed walk on $K_7$ that covers each of
the $21$ edges exactly once is precisely an Eulerian circuit,
giving
\begin{equation}
  \mathcal{A}[\,e \to \sigma\,]
  \;\sim\;
  \prod_{\{i, j\} \in E(K_7)} \sqrt{\alpha}
  \;=\; \alpha^{21/2} .
\end{equation}
The $\alpha$-\emph{exponent} $21/2$ is thus determined
structurally by:
(i) $21 = \dim \mathfrak{so}(7)$ via the $\mathrm{Cl}(0, 7)$
construction;
(ii) $PSL(2, 7)$-equivariance of the physical amplitude, forcing
one coupling per $\mathfrak{so}(7)$ generator;
(iii) Paper~13's established $\sqrt{\alpha}$-per-crossing rule on
open transition paths.

What remains conjectural is the functional \emph{form} of the
overall amplitude (specifically: that the $PSL(2, 7)$-equivariant
leading-order amplitude is a product over all $21$ generators, as
opposed to a sum or a different structural combination) and the
values of the dimensionless prefactor and NLO correction.

We regard the current status as intermediate between Balmer's
formula (1885) and the full Bohr model: the structural anatomy
is complete, the $\alpha$-exponent derivation is rigorous modulo
the amplitude-form assumption, and only the dimensionless
normalisations remain for a full dynamical derivation.

\subsection{Relation to $\mathfrak{spin}(7)$}

The integers $\CA = 7$, $\CA + 1 = 8$, and $21 = \CA(\CA-1)/2$
correspond exactly to the \textbf{vector, spinor, and adjoint}
representation dimensions of $\mathrm{Spin}(7)$.  The hierarchy
formula reads
\begin{equation}
  \frac{m_e}{m_\text{Pl}}
  = \frac{\dim(\text{spinor})}{\dim(\text{vector})}
    \left(1 + \frac{\alpha}{\dim(\text{vector})}\right)
    \alpha^{\dim(\text{adjoint})/2},
\end{equation}
with all three dimensions arising from the same Lie group.

This structural identification has a concrete \textbf{matrix
realisation} via the octonion Clifford algebra $\mathrm{Cl}(0,7)$.
Despite octonion non-associativity, the left-multiplication
operators $L_i := L_{e_i}$ for $i = 1, \ldots, 7$ satisfy the
Clifford relations
\begin{equation}
  \{L_i, L_j\} \;=\; L_i L_j + L_j L_i \;=\; -2\,\delta_{ij}\,I_8
\end{equation}
\emph{exactly}: because octonions are an alternative algebra, the
associator $[a, b, c] := (ab)c - a(bc)$ is antisymmetric in its
first two arguments, so the associator terms in the anticommutator
cancel identically.  Consequently $L_i$ are $\mathrm{Cl}(0, 7)$
gamma matrices acting on $\mathbb{R}^8 = \mathbb{O}$, and the
bivectors
\begin{equation}
  \mathbf{i} := L_1 L_2,\qquad
  \mathbf{j} := L_1 L_3,\qquad
  \mathbf{k} := L_2 L_3
\end{equation}
satisfy the quaternion relations $\mathbf{i}^2 = \mathbf{j}^2 =
\mathbf{k}^2 = -I$, $\mathbf{i}\mathbf{j} = \mathbf{k}$, etc.,
by direct computation.  This gives $\mathrm{Spin}(3) \subset
\mathrm{Spin}(7)$.  The binary tetrahedral group $2T \subset
SU(2) = \mathrm{Spin}(3)$ then embeds faithfully into
$\mathrm{Spin}(7)$ via
\begin{equation}
  M(q) \;=\; a\, I + b\,(L_1 L_2) + c\,(L_1 L_3) + d\,(L_2 L_3),
  \qquad q = a + bi + cj + dk \in 2T.
\end{equation}
This is verified as a genuine $2T \to \mathrm{Spin}(7) \subset SO(8)$
embedding: all $24$ matrices are in $SO(8)$, $M(-e) = -I_8$ (center
preserved), the group law $M(q_1 q_2) = M(q_1) M(q_2)$ holds exactly
for all $576$ pairs, and a $\mathrm{Spin}(7)$-invariant $4$-form is
preserved by all $24$ elements (see supplementary materials).

Thus Paper~15's dimensional chain $(7, 8, 21)$ has an explicit
matrix realisation: $2T$ acts on the $8$-dim spinor via $M$, the
$7$-dim vector is $\operatorname{Im}(\mathbb{O})$, and the $21$-dim
adjoint $\mathfrak{so}(7)$ organises the coupling channels.

The earlier identification $168 = |PSL(2,7)|$ retains interest as a
secondary coincidence --- $PSL(2,7)$ acts $2$-transitively on the
Fano plane with incidence count $21$, and $K_7$ on the Heegaard
torus of $S^3/2I$ carries a natural $PSL(2,7)$ action --- but the
primary structural source of Paper~15's integers is $2T$ and its
McKay / $\mathrm{Cl}(0,7)$ / $\mathrm{Spin}(7)$ chain.

\subsection{The Kerr connection}

The electron's Kerr--Newman ring radius is
$a = \hbar/(2m_e c) = \xi_\text{SM}/2$---exactly half the BPS
healing length.  Carter's $g = 2$~\cite{Carter1968} for the
Kerr metric matches Dirac's quantum result, and the charge
radius $r_Q = \sqrt{\alpha}\,\ell_\text{Pl}$ is the geometric
mean of topology and quantum gravity.  The BPS vortex core
regularises what would otherwise be a Kerr ring singularity.

\subsection{The residual}

The NLO prediction matches CODATA at $0.013\%$.  The NNLO
correction $(1 + \alpha/7 + \alpha^2/49 + \cdots) = 1/(1 -
\alpha/7)$ shifts $G$ by only $\alpha^2/49 \approx 10^{-6}$,
well below the current residual.  The remaining $0.013\%$ likely
reflects the running of $\alpha$ between the Thomson limit and
the gravitational-crossover scale, or experimental uncertainty
in $G$ itself (the CODATA central value has shifted by
$\sim 0.04\%$ between adjustments).


%=====================================================================
\section{Lagrangian context: the result as $\mathcal{L}_5$ of
the minimal three-field NWT theory}
\label{sec:lagrangian-context}
%=====================================================================

The gravitational hierarchy examined in this paper fits as one
slice of a larger structural picture developed in a companion
manuscript.  A proposed minimal relativistic Lagrangian whose
torus-knot solitons are intended to reproduce NWT's observational
footprint is a three-field theory on $\mathbb{R}^{1,3}$ (or, in
the GUT-extended UV completion, on a Kaluza--Klein space with
internal $S^3/2I$):

\begin{equation}
\label{eq:LNWT}
\mathcal{L}_\text{NWT}
  =  |D_\mu \psi|^2  -  \tfrac{1}{4} F_{\mu\nu} F^{\mu\nu}
  -  \tfrac{\lambda}{4}\bigl(|\psi|^2 - v^2\bigr)^2
  +  \tfrac{f_\pi^2}{2}\,(\partial_\mu n^a)^2
  -  \tfrac{1}{4 e_\text{Sk}^2}
     (\partial_\mu n^a \partial_\nu n^b - \partial_\nu n^a
        \partial_\mu n^b)^2
  +  \theta\,Q_H[n] .
\end{equation}

The fields are a complex scalar condensate $\psi$, an abelian
gauge field $A_\mu$, and an $S^2$-valued Faddeev unit field
$n^a$ (which can be written as $n^a = \psi^\dagger \sigma^a
\psi/(\psi^\dagger \psi)$ in the $CP^1$ extension).  The
couplings $\lambda$, $f_\pi$, $e_\text{Sk}$, $\theta$ are fixed
at the BPS critical point $\lambda = e^2/2$, the Derrick
equilibrium $\kappa = R/a_0 = \pi^2$, and the CPT-invariant
$\theta = 0$ or $\pi$.  No free fit parameters beyond the
overall mass scale $v = m_e$.

The full structural programme decomposes into five layers,
labelled $\mathcal{L}_1$ through $\mathcal{L}_5$, of which the
present paper is the fifth:

\begin{itemize}
  \item $\mathcal{L}_1$: minimal field content $(\psi, A_\mu,
        n^a)$, with the $\mathrm{Spin}(7)/\mathrm{Cl}(0,7)/2T$
        structure of \S\ref{sec:mechanism}--\ref{sec:hierarchy}
        emerging on the moduli space of torus-knot solitons,
        not as a fundamental gauge symmetry.
  \item $\mathcal{L}_2$: BPS abelian Higgs sector
        ($|D\psi|^2 + F^2 + V$) at $\lambda = e^2/2$.
        Bogomolny bound saturated; line tension
        $\mu_\text{BPS} = 2\pi v^2$ verified to $\sim 10^{-3}\%$.
  \item $\mathcal{L}_3$: Skyrme--Faddeev quartic + Hopf
        $\theta$-term, ensuring stable finite-size knotted
        solitons via Derrick scaling and locking the
        topological sector to $Q_H = p \cdot m$ for a
        $T(p, q)$ torus knot with phase winding $m$.
  \item $\mathcal{L}_4$: Paper~6 mass formula for $24$ hadrons
        and leptons ($1.06\%$ median deviation, no fit parameters
        beyond the anchor $m_e$) is obtained from $\mathcal{L}_2$'s
        line tension evaluated on $\mathcal{L}_3$'s torus-knot
        solitons via the Kelvin--Saffman thin-ring approximation.
  \item $\mathcal{L}_5$: \emph{this paper}.  The gravitational
        hierarchy
        $m_e/m_\text{Pl} = (8/7)(1 + \alpha/7)\alpha^{21/2}$
        is interpreted as arising from the one-loop Casimir of
        $\mathcal{L}_2$ on $S^3/2I$, with the $T(2,3)$ trefoil
        sector selected by $\mathcal{L}_3$'s Hopf $\theta$-term.
\end{itemize}

Each factor in the gravitational formula can be associated with a
proposed origin in Eq.~\ref{eq:LNWT}.
Table~\ref{tab:lagrangian-attribution} summarises the mapping;
the attributions are best read as a coherent map of the model's
internal logic rather than as a completed microscopic derivation
(for which see \S\ref{sec:mechanism} on what is rigorous versus
structurally grounded).

\begin{table}[H]
\centering
\renewcommand{\arraystretch}{1.25}
\begin{tabular}{lll}
\toprule
\textbf{Factor} & \textbf{Lagrangian origin} & \textbf{Reference} \\
\midrule
$8/7$ &
  $\mathcal{L}_1$ field content: &
  \S\ref{sec:mechanism}, \S\ref{sec:hierarchy} \\
& $\dim(\mathrm{Spin}(7)\text{ spinor}) /
   \dim(\mathrm{Spin}(7)\text{ vector})$ & \\
& matter / gauge propagator ratio & \\
\midrule
$\alpha^{21/2}$ &
  $\mathcal{L}_2$ gauge sector: &
  \S\ref{sec:mechanism} \\
& product over $21$ edges of Eulerian & \\
& circuit on $K_7$, each contributing $\sqrt{\alpha}$ & \\
\midrule
$(1 + \alpha/7)$ &
  $\mathcal{L}_2$ one-loop NLO: &
  \S\ref{sec:nlo} \\
& seven $K_7$ vertex self-energies, & \\
& each of order $\alpha/49$ & \\
\midrule
$\theta\,Q_H$ lock &
  $\mathcal{L}_3$ Hopf $\theta$-term: &
  \S\ref{sec:mechanism} \\
& selects the $(2,3)$ trefoil sector & \\
& via the $2I$ surgery equivalence & \\
\bottomrule
\end{tabular}
\caption{Attribution of each factor in
$m_e/m_\text{Pl} = (8/7)(1 + \alpha/7)\alpha^{21/2}$ to a
specific term in the minimal NWT Lagrangian
(Eq.~\ref{eq:LNWT}).  Companion Paper~16~\cite{paper16}
develops the full $\mathcal{L}_1$--$\mathcal{L}_5$ chain.}
\label{tab:lagrangian-attribution}
\end{table}

The forthcoming Paper~16~\cite{paper16} assembles
$\mathcal{L}_1$--$\mathcal{L}_5$ into a single coherent
treatment, including the $SO(10)$ UV completion at
$E_\text{GUT} \approx 7.4 \times 10^{15}$~GeV (which adds
$33$ heavy gauge bosons orthogonal to the present paper's
emergent $\mathrm{Spin}(7)$ structure), and three concrete
falsifiers: Hyper-Kamiokande proton-decay reach
$\tau(p \to e^+ \pi^0) \sim 1.4 \times 10^{35}$~yr,
non-MSSM coupling unification at $\alpha_\text{GUT} = 1/40$,
and a GUT phase-transition gravitational-wave signature near
$7.4$~GHz.  The present paper provides $\mathcal{L}_5$.

%=====================================================================
\section{Conclusion}
\label{sec:conclusion}
%=====================================================================

Within the accounting developed here, the trefoil knot organises
a substantial portion of the framework's output:
\begin{enumerate}\setlength\itemsep{2pt}
  \item Its \textbf{crossings} generate $\su(3)$ and the fine
    structure constant $\alpha$ (Papers~8, 11, 13--14).
  \item Its \textbf{Dehn surgery} produces the Poincar\'e sphere
    $S^3/2I$ and the icosahedral group $A_5$, giving nuclear
    saturation $k_\text{sat} = 5$ (Paper~17) and the cinquefoil
    eigenvalue $q^2 = 25$ in the $\alpha$ formula (Paper~14).
  \item Its \textbf{Poincar\'e spectrum} $\lambda_1 = 168 =
    |\mathrm{Irr}(2T)| \times |2T| = 7 \times 24$ (McKay
    correspondence, extended $E_6$ Dynkin) supplies the
    integers entering the gravitational hierarchy
    $m_e/m_\text{Pl} = (8/7)(1{+}\alpha/7)\,\alpha^{21/2}$
    (this paper).
  \item Its \textbf{$\mathrm{Spin}(7)$ embedding} $2T \subset
    \mathrm{Spin}(7)$ via the octonion Clifford algebra
    $\mathrm{Cl}(0, 7)$ provides an explicit matrix realisation;
    the three integers $(7, 8, 21)$ in the gravitational hierarchy
    formula are the vector, spinor, and adjoint dimensions of
    $\mathrm{Spin}(7) = B_3$ (this paper, Sec.~\ref{sec:mechanism}).
\end{enumerate}

The predicted gravitational constant agrees numerically with
CODATA at the $0.013\%$ level.  Within the dimensionless content
of the framework, no fit parameter remains (modulo the
structural-vs-dynamical caveats of \S\ref{sec:mechanism} and
\S\ref{sec:zero}): the one-knot accounting reduces the Standard
Model plus gravity from $\sim 21$ dimensionless parameters to
zero.

\section*{Acknowledgements}

The spectral connection $\lambda_1 = 168 = 8 \times 21$ was
discovered during an exploratory session on gravitational
self-consistency and Kerr naked singularities, following the
publication of Paper~14.  The NLO pattern $(1 + \alpha/\CA)$
was identified by systematic stress-testing of the leading-order
formula against alternative integer decompositions.

We thank four AI reviewers---Gemini (Google), Perplexity,
ChatGPT (OpenAI), and Claude Mobile (Anthropic, Opus~4.7)---for
critical reviews of the companion Paper~14 that sharpened the
uniqueness arguments, statistical framework, and braid-algebra
presentation carried forward into this work.

All analysis code, simulation scripts (BPS profile solver,
Poincar\'e eigenvalue computation, holonomy integrator,
crossing algebra, knot selection scan), and the \LaTeX\ source
for this manuscript are available at
\url{https://github.com/JimGalasyn/null-worldtube}.

%=====================================================================
\begin{thebibliography}{99}
%=====================================================================

\bibitem{paper13}
  J.~P.~Galasyn and C.~Th\'eodore,
  ``The Standard Model from NWT Topology (Paper~13),''
  Zenodo (2026), DOI:~10.5281/zenodo.19635239.

\bibitem{paper14}
  J.~P.~Galasyn and C.~Th\'eodore,
  ``Integers as Output (Paper~14),''
  Zenodo (2026), DOI:~10.5281/zenodo.19654507.

\bibitem{paper16}
  J.~P.~Galasyn and C.~Th\'eodore,
  ``The NWT Lagrangian: a Three-Field Theory of Particles and
  Gravity (Paper~16),''
  in preparation (2026).

\bibitem{paper17}
  J.~P.~Galasyn and C.~Th\'eodore,
  ``Nuclear Binding from Topological Linking (Paper~17),''
  in preparation (2026).

\bibitem{Carter1968}
  B.~Carter,
  ``Global structure of the Kerr family of gravitational fields,''
  \emph{Phys.\ Rev.}\ \textbf{174}, 1559 (1968).

\bibitem{Rolfsen}
  D.~Rolfsen,
  \emph{Knots and Links},
  AMS Chelsea Publishing (2003).

\bibitem{AlonsoIzquierdo2016}
  A.~Alonso-Izquierdo, W.~Garc\'ia Fuertes, and J.~Mateos Guilarte,
  ``One-loop mass shift formula for kinks and
  self-dual vortices,''
  \emph{Phys.\ Rev.}\ D \textbf{94}, 125003 (2016).

\bibitem{BaackeLavrelashvili2008}
  J.~Baacke and N.~Kevlishvili,
  ``One-loop corrections to the quantum mass of the
  abelian Higgs vortex,''
  \emph{Phys.\ Rev.}\ D \textbf{78}, 085008 (2008).

\bibitem{public-repo}
  J.~P.~Galasyn,
  ``Null Worldtube Theory: public code repository,''
  \url{https://github.com/JimGalasyn/null-worldtube}.

\end{thebibliography}

\end{document}
